\section{Conclusion}

In this paper, we propose PRIME, a unified \textbf{training-free} ANN-to-SNN conversion framework that tackles both information redundancy in time-step extension and computational redundancy in auxiliary neuronal dynamics.
%Unlike training-dependent methods, our method converts pretrained ANN weights in a single forward pass for training-free conversion.
At the encoding level, the proposed MTNs with RMPA tighten the error bound from $\theta/2T$ of rate coding to $\theta/2^{T \cdot D+1}$, enabling near-lossless conversion with $T \leq 2$; 
at the inference level, thresholds and values are read directly from FP binary representation, eliminating membrane potential accumulation and multi-threshold comparison.
%Under $T = 2$, $D = 3$, PRIME attains 10.5\% of FP32 multiplication energy (35.5\% of FP16) and 40.6\% of FP32 bit-transmission cost (81.3\% of FP16).
These results establish a hardware-aligned design principle for efficient ANN-to-SNN conversion.



%amplify the residual membrane potential at each iteration, tightening the error bound from the linear $\theta/2^{T}$ of rate coding to an exponential $\theta/2^{T \cdot D+1}$ and enabling near-lossless conversion at $T \leq 2$; at the inference level, thresholds and values are read directly from the floating-point binary representation, eliminating membrane potential accumulation and multi-threshold comparison.
%Under $T = 2$, $D = 3$, PRIME attains 10.5\% of FP32 multiplication energy (35.5\% of FP16) and 40.6\% of FP32 bit-transmission cost (81.3\% of FP16).
%These results establish a hardware-aligned design principle for efficient ANN-to-SNN conversion.
% In this paper, we propose PRIME, a unified \textbf{training-free} ANN-to-SNN conversion framework that jointly addresses information redundancy in time-step extension and computational redundancy in auxiliary neuronal dynamics. Unlike training-dependent methods~\cite{you2024spikezip,xing2024spikellm}, PRIME requires no retraining or fine-tuning---all conversions are performed in a single forward pass on the pretrained ANN weights, making it applicable to proprietary or frozen models.
% At the encoding level, integration-coding-based MTNs leverage the RMPA mechanism to actively amplify residual membrane potential at each iteration, achieving an exponentially tighter error bound of $\theta/2^{T \cdot D+1}$ versus the linear bound $\theta/2^{T}$ of rate coding, enabling near-lossless conversion at $T \leq 2$. At the inference level, spike thresholds and values are read directly from the floating-point binary representation, eliminating membrane potential accumulation and multi-threshold comparison; the resulting spike-driven multiplication requires at most a single adder per cycle when $D \leq 3$. Under $T = 2$, $D = 3$, energy consumption reaches 10.5\% of FP32 multiplication (35.5\% of FP16) and bit-transmission cost reaches 40.6\% of FP32 (81.3\% of FP16).
%These figures rely on analytical estimation; however, all underlying primitives (bit-shift, integer addition, and leading-one detection) are natively supported by standard digital logic, making physical implementation a natural next step.